% Ч1.4 — Увеличение и уменьшение на p% (множитель 1 ± p/100)
% Старая методичка: стр. 152 (старое задание 16)

\begin{tema}{Ч1.4}
\opr{Увеличить} число $x$ на $p\%$ — значит умножить его на $1+\frac{p}{100}$:
\[
  x+x\cdot\frac{p}{100}=\left(1+\frac{p}{100}\right)x.
\]
\opr{Уменьшить} число $x$ на $p\%$ — значит умножить его на $1-\frac{p}{100}$:
\[
  x-x\cdot\frac{p}{100}=\left(1-\frac{p}{100}\right)x.
\]
Число $k=1\pm\frac{p}{100}$ называют \opr{повышающим} (понижающим) коэффициентом.

\begin{primer}
$+20\%$ — умножить на $1{,}2$; \ $-15\%$ — умножить на $0{,}85$;
\ $+12{,}5\%$ — умножить на $1{,}125=\frac98$.
\end{primer}

На этом построена экономическая задача (№\,13): каждый год долг или вклад
умножается на $k$ — см. \tref{М2.1}.
\end{tema}
